The empirical results established that the hybrid HMM framework generates synthetic equity paths whose distributional and temporal properties closely match those of historical data, both in-sample and out-of-sample. In this section we interpret the key findings, discuss their practical implications, and identify the limitations that inform the scope of these conclusions.

% ── Figure 7: Multi-Asset SIM Extension ──────────────────────────────────────
\begin{figure}[tp]
    \centering
    \includegraphics[width=\textwidth]{Fig7-Multi-Asset-SIM.pdf}
    \caption{Multi-asset extension via the Single-Index Model (SIM) across 424 S\&P~500
    constituents.  Panel~(a) summarizes the SIM regression fit quality by GICS sector.
    Panel~(b) shows that KS distributional consistency is maintained across the full asset
    universe when paths are generated via the factor structure
    $\hat{G}_{i,t} = \hat{\alpha}_i + \hat{\beta}_i G_{{\rm SPY},t} + \hat{\eta}_{i,t}$.
    Panel~(c) illustrates representative marginal density comparisons for three assets spanning
    a wide range of systematic risk exposure.}
    \label{fig:multi_asset}
\end{figure}


\subsection{The role of the jump-duration mechanism}
The comparison between HMM-NJ and HMM-WJ isolated the contribution of the Poisson jump-duration mechanism. Without jumps, the Markov chain transitioned out of tail states within one to two steps on average, producing an ACF of $|G_t|$ that decayed too quickly and failing to reproduce the volatility clustering observed in the data. Introducing jumps forced the model to dwell in high-volatility regimes for empirically realistic durations, and approximately 24\% of simulated paths contained at least one such episode, sufficient to shift the ensemble-level ACF into close agreement with the empirical profile. This result is consistent with the broader literature on regime-switching models, where standard HMM specifications have repeatedly been found to undergenerate volatility persistence \cite{ryden_stylized_1998, bulla_stylized_2006}, and it demonstrated that a minimal two-parameter extension ($\epsilon$, $\lambda$) could resolve the issue without the computational overhead of hierarchical or semi-Markov alternatives. The cross-asset analysis (Online Appendix~S5) reinforced this interpretation: the grid search selected $\lambda^{*}=160$ for NVDA and JPM (assets with pronounced volatility clustering) but only $\lambda^{*}=30$ for JNJ (a defensive stock with shorter-lived volatility episodes), indicating that $\lambda$ adapted automatically to each asset's clustering intensity. For JNJ, the jump mechanism contributed almost no improvement over the standard HMM (ACF-MAE 0.027$\to$0.026), consistent with the expectation that low-volatility assets exhibit less persistent tail-state occupancy and therefore require shorter, or fewer, jump episodes.

The tradeoff between distributional fidelity and temporal structure, visible in the GARCH comparison, deserves emphasis. GARCH(1,1) achieved the lowest ACF-MAE because its variance equation was explicitly parameterized to match autocorrelation dynamics, yet its marginal distribution was systematically inconsistent with the data (IS KS pass rate of 5.5\%). The hybrid HMM reversed this tradeoff: the Laplace-mixture emissions delivered high distributional fidelity, and the jump mechanism provided sufficient temporal structure. Neither model dominated on all metrics simultaneously, but the HMM framework provided a more balanced quality profile for synthetic data applications where both distributional and temporal fidelity matter. Crucially, the volatility-clustering strength of HMM-WJ is not fixed: increasing $\epsilon$ raises the fraction of paths that contain jump episodes and drives ACF-MAE toward zero, while decreasing $\epsilon$ recovers the pure Markov baseline. The grid search selected the value of $\epsilon$ that jointly optimized temporal and distributional fidelity, so the observed ACF-MAE reflects the best achievable balance for SPY rather than a structural limitation of the mechanism.

\subsection{Practical implications}
The proposed framework offers practical utility for both risk management and the broader synthetic data community. Generative models that reproduce the stylized facts of financial markets enable stress testing under market regimes that are statistically plausible but not historically observed, addressing a key limitation of scenario libraries derived solely from empirical records. The interpretable regime structure, where each hidden state corresponds to a quantile-defined segment of the excess growth rate distribution, facilitates communication between quantitative analysts and risk managers; states can be labeled (e.g., \textit{crash}, \textit{bear}, \textit{neutral}, \textit{bull}, \textit{rally}) and linked to economic narratives. The model's computational efficiency further supports Monte Carlo-based risk metrics such as Value-at-Risk and Conditional Value-at-Risk at scale across large asset universes \cite{glasserman_monte_2003}.

The avoidance of the Baum-Welch EM algorithm, made possible by the quantile-based state partition that explicitly assigns each observation to a hidden state, eliminates convergence sensitivity to initialization and reduces computational cost. This property matters for large-scale pipelines: fitting the model to 424 assets requires only empirical counting and a two-parameter grid search, making it feasible to regenerate synthetic scenarios on a daily or weekly cadence as new data arrive \cite{ohara_market_1998, manganelli_2004}. More broadly, the quality evaluation methodology demonstrated here, combining distributional fidelity tests (KS, AD) with temporal structure metrics (ACF-MAE) and reporting standard errors across Monte Carlo ensembles, provides a reusable template for assessing synthetic time series quality in other domains where temporal dependence structures are critical to downstream applications \cite{bauer_jdiq_ts_synthesis_2024}.

\subsection{Limitations}
Several limitations inform the interpretation of these results. First, the model assumes stationarity of the transition matrix and jump hyperparameters across the full in-sample window. The out-of-sample ACF analysis (Figure~\ref{fig:statistical_validation}, panel~d) provided direct evidence of this limitation: the IS-calibrated $(\epsilon^*,\lambda^*)$ underestimated volatility clustering persistence in the 2025 test window, consistent with the elevated macro uncertainty of that period. Structural breaks in market microstructure, such as those accompanying regulatory changes or the COVID-19 shock, may not be adequately captured by a single static transition matrix.

Second, the quantile-based state partition is fixed at estimation time and does not adapt to evolving volatility regimes. The boundaries are calibrated on the full in-sample distribution and may overweight the central mass relative to the tails during periods of prolonged market stress.

Third, the Single-Index Model extension imposes a linear, single-factor model of cross-sectional dependence. The median IS KS pass rate of 66.7\% across 424 assets confirmed that this decomposition could not capture the heterogeneous tail behavior, skewness, and sector-specific dynamics present in individual equities \cite{sharpe_simplified_1963, fama_french_1993}. However, the architecture is modular: the HMM-WJ generative engine for the market factor is decoupled from the factor model used to propagate paths, and replacing the SIM with a richer specification would not require modifying the underlying Markov chain machinery.

Finally, the out-of-sample evaluation covers a single 249-day window (full calendar year 2025), which limits the statistical power of out-of-sample inference. A longer or rolling evaluation would provide a more robust assessment of generalization performance.
